\documentclass[journal,10pt]{IEEEtran}

\usepackage[utf8]{inputenc}
\usepackage[spanish,es-nodecimaldot]{babel}
\usepackage{amsmath,amssymb,amsfonts,bm}
\usepackage{algorithmic}
\usepackage{graphicx}
\usepackage{textcomp}
\usepackage{xcolor}
\usepackage{booktabs}
\usepackage{multirow}
\usepackage{hyperref}
\usepackage{cite}
\usepackage{listings}
\usepackage{microtype}

\hypersetup{
    colorlinks=true,
    linkcolor=blue!80!black,
    filecolor=magenta,      
    urlcolor=cyan!70!black,
    citecolor=red!70!black
}

\lstset{
  basicstyle=\ttfamily\scriptsize,
  breaklines=true,
  frame=single,
  backgroundcolor=\color{gray!8},
  keywordstyle=\color{blue}\bfseries,
  commentstyle=\color{green!50!black}\itshape,
  stringstyle=\color{red!70!black}
}

\begin{document}

\title{Convergencia de Seguridad Multinivel (MLS), Notarización en Blockchain e Inferencia Asistida por IA On-Premise para Ecosistemas Documentales Críticos}


\author{
  \IEEEauthorblockN{David Enrique Duarte Gonzalez}\\
  \IEEEauthorblockA{
    \textit{Universidad de Carabobo}\\
    \textit{Dirección de Informática}\\
    Valencia, Venezuela\\
    \texttt{\{deduarte1\}@uc.edu.ve}
  }
}

\maketitle

\begin{abstract}
La administración de documentos y comunicaciones oficiales en instituciones públicas y sectores de seguridad y defensa nacional demanda garantías intransigentes de confidencialidad jerárquica, autenticidad comprobable, integridad temporal y no-repudio estricto. La adopción indiscriminada de soluciones comerciales de gestión documental basadas en computación en la nube pública (\textit{Cloud SaaS}) expone a los organismos a vulnerabilidades críticas de exfiltración foránea, pérdida de soberanía sobre los datos, fugas en reposo y alteración retroactiva de correlativos institucionales. En este artículo, se presenta la concepción, fundamentación matemática, diseño arquitectónico e implementación empírica de \textbf{Polibio}: una plataforma institucional de microservicios contenerizados diseñada bajo la doctrina de soberanía digital y confianza cero (\textit{Zero Trust}). La solución unifica de manera sinérgica cuatro pilares defensivos: (1) un modelo de Seguridad Multinivel (MLS) conforme a los teoremas de Bell-LaPadula integrado con políticas a nivel de fila (\textit{Row-Level Security}, RLS) en bases de datos relacionales; (2) un esquema criptográfico de sobre híbrido autenticado (AES-256-GCM y RSA-OAEP) con descifrado exclusivo al vuelo en memoria RAM que erradica por diseño la presencia de texto plano en discos secundarios; (3) un libro mayor inmutable basado en Blockchain que notariza documentalmente las firmas de una Autoridad Certificadora Raíz (PKI) soberana e introduce el registro forense inmutable de intentos de acceso denegados; y (4) un motor de inferencia neuronal on-premise aislado que provee asistencia de redacción inteligente con cero telemetría externa. La evaluación experimental confirma una sobrecarga criptográfica sub-milisegundo, resistencia total a la falsificación de correlativos y la preservación matemáticamente verificable del principio de no-repudio.
\end{abstract}

\begin{IEEEkeywords}
Seguridad Multinivel (MLS), Bell-LaPadula, Cifrado de Sobre AES-256-GCM, Blockchain Permisionado, Notarización Inmutable, Trazabilidad Forense, Row-Level Security, Soberanía Digital.
\end{IEEEkeywords}

\section{Introducción}
\IEEEPARstart{L}{a} transformación digital de los aparatos de Estado y las fuerzas de seguridad demanda una revisión radical de las premisas bajo las cuales se conciben los sistemas de trámite y gestión de documentos oficiales. En el marco de la ciberguerra asimétrica y la soberanía tecnológica, los flujos documentales gubernamentales no representan meros repositorios de texto, sino los vehículos vinculantes a través de los cuales se formalizan resoluciones ejecutivas, instrucciones operativas, directivas de inteligencia y asignaciones de recursos institucionales. 

No obstante, los sistemas de gestión documental contemporáneos adolecen de severas limitaciones arquitectónicas:
\begin{enumerate}
    \item \textbf{Dependencia Externa y Fuga de Soberanía:} La proliferación de plataformas propietarias en esquemas de nube pública extranjera expone la información institucional a marcos jurídicos transfronterizos y puertas traseras de actores hegemónicos \cite{nist_sp800_207_zerotrust}.
    \item \textbf{Vulnerabilidad de los Datos en Reposo (\textit{At-Rest Leakage}):} Los sistemas tradicionales dependen de cifrado de volumen a nivel de disco (\textit{Full Disk Encryption}). Si bien esto protege contra el robo físico de un disco duro, resulta completamente ineficaz ante ataques internos o escalamiento de privilegios en el sistema operativo anfitrión, donde el motor de base de datos mantiene las tablas en texto claro.
    \item \textbf{Maleabilidad y Manipulación Retroactiva de Correlativos:} Las bases de datos relacionales tradicionales permiten que un administrador malicioso modifique números de oficios, reasigne fechas de emisión o elimine registros de auditoría sin dejar huellas forenses indelebles.
    \item \textbf{Exfiltración Clandestina a través de Asistentes IA:} La integración apresurada de modelos de lenguaje basados en APIs comerciales de terceros introduce un vector de fuga masivo, donde el contenido de documentos clasificados es transmitido a servidores remotos para su procesamiento \cite{carlini2021extracting}.
\end{enumerate}

Para resolver estas problemáticas estructurales, en este artículo introducimos \textbf{Polibio}, un ecosistema integral on-premise concebido para operar bajo condiciones de aislamiento de red (\textit{air-gap} o red soberana restringida). El sistema materializa la convergencia sinérgica entre modelos formales de control de acceso multinivel, criptografía de última generación y libros mayores inmutables.

\subsection{Contribuciones Principales}
Las contribuciones fundamentales de esta investigación se sintetizan en los siguientes aspectos:
\begin{itemize}
    \item \textbf{Formalización e Implementación de MLS Dual:} Se formaliza e implementa un modelo Bell-LaPadula jerárquico de 5 niveles que opera tanto en la capa perimetral de microservicios mediante encabezados criptográficos de sesión, como en la capa física de almacenamiento mediante políticas de seguridad a nivel de fila (\textit{Row-Level Security}, RLS) en PostgreSQL.
    \item \textbf{Cifrado de Sobre Híbrido con Cero Fuga en Reposo:} Se introduce un protocolo de sobre criptográfico autenticado que combina claves efímeras simétricas AES-256-GCM y envoltura asimétrica RSA-OAEP. Los documentos de nivel reservado o superior se almacenan saneados en disco secundario mediante marcadores seguros, garantizando que el descifrado suceda exclusivamente en memoria volátil transitoria (\textit{in-flight}) ante credenciales criptográficas válidas.
    \item \textbf{Libro Mayor Blockchain Forense con Registro de Infracciones:} A diferencia de los libros mayores convencionales que solo indexan documentos exitosos, Polibio registra en su cadena de bloques inmutable cada rechazo de acceso por violación de nivel MLS. Esto transforma la cadena en un registro forense vinculante e inalterable ante auditorías de contrainteligencia.
    \item \textbf{Asistencia Neuronal Soberana On-Premise:} Se integra un motor de inferencia de modelos abiertos compactos ejecutado localmente en contenedor sin comunicación con el exterior, blindando los borradores de cualquier vector de exfiltración.
\end{itemize}

\subsection{Estructura del Artículo}
El resto de este documento se organiza de la siguiente manera: la Sección II examina los trabajos previos y el marco teórico; la Sección III detalla la arquitectura global de microservicios; la Sección IV profundiza en el modelo MLS y el cifrado de sobre; la Sección V presenta la notarización en Blockchain y la auditoría forense; la Sección VI describe la inferencia neuronal soberana; la Sección VII expone la evaluación experimental; y finalmente, la Sección VIII concluye el trabajo delineando futuras líneas de investigación.

\section{Marco Teórico y Trabajos Relacionados}
\subsection{Modelos de Seguridad Multinivel y Teoremas de Bell-LaPadula}
El paradigma de Seguridad Multinivel (MLS) propuesto seminalmente por Bell y LaPadula \cite{bell1976secure} y formalizado algebraicamente por Sandhu \cite{sandhu1993lattice} fundamenta el control de acceso mandatorio (MAC) para entidades militares y gubernamentales. En dicho marco, los sujetos y los objetos son asignados a un reticulado de seguridad $\mathcal{L} = (L, \le)$ donde se imponen dos axiomas esenciales:
\begin{enumerate}
    \item \textbf{Propiedad de Seguridad Simple (No-Read-Up):} Un sujeto $S$ con habilitación $f_S$ puede leer un objeto $O$ con clasificación $f_O$ si y solo si $f_O \le f_S$.
    \item \textbf{Propiedad Estrella ($\star$-Property o No-Write-Down):} Un sujeto $S$ no puede escribir o rebajar información clasificada a un objeto de menor nivel jerárquico ($f_S \le f_O$).
\end{enumerate}
La mayoría de las implementaciones modernas restringen MLS al sistema operativo (e.g., SELinux). Polibio extiende este modelo hacia la capa de microservicios y bases de datos relacionales en tiempo real.

\subsection{Criptografía de Sobre y Cifrado Autenticado con Datos Asociados (AEAD)}
El modo Galois/Counter Mode (GCM) de AES \cite{nist_sp800_38d_gcm} provee confidencialidad semántica y verificación de integridad criptográfica en una única pasada computacional, produciendo una etiqueta de autenticación (\textit{tag}) $T \in \{0, 1\}^{128}$. El esquema de sobre híbrido (\textit{Envelope Encryption}) utiliza AES-256-GCM para el cuerpo del mensaje acoplado con derivación asimétrica RSA-OAEP \cite{rfc8017_rsa_oaep}, mitigando vulnerabilidades contra ataques de texto cifrado adaptativo (\textit{IND-CCA2}) \cite{boneh2023graduate}.

\subsection{Tecnologías de Libro Mayor Distribuido y Notarización Temporal}
Los protocolos de sellado de tiempo de Haber y Stornetta \cite{haber1991how} demostraron la invulnerabilidad del encadenamiento criptográfico con funciones de dispersión unidireccional (SHA-256). Plataformas modernas de cadena de bloques permisionadas \cite{androulaki2018hyperledger, xu2019architecture} permiten establecer redes de consenso cerrado donde se preserva la cronología de transacciones con valor probatorio en auditorías externas \cite{gaetani2019blockchain}.

\subsection{Privacidad y Riesgos de Exfiltración en Asistentes Inteligentes}
Las investigaciones de Carlini et al. \cite{carlini2021extracting} revelan que los Modelos de Lenguaje de Gran Escala (LLM) memorizan fragmentos verbatim de su contexto de inferencia. Cuando una institución pública utiliza servicios remotos comerciales, la información gubernamental confidencial queda sujeta a fuga a través de telemetría o retención de prompts. La ejecución de modelos abiertos como Qwen2 \cite{yang2024qwen2} en infraestructura interna elimina radicalmente dicho riesgo.

\section{Arquitectura del Ecosistema Polibio}
\subsection{Topología de Malla de Microservicios Contenerizados}
Polibio se estructura como una constelación de seis microservicios desacoplados orquestados mediante contenedores Docker, comunicados mediante una red mallada aislada (\texttt{polibio-mesh}) y gobernados por patrones arquitectónicos de microservicios seguros \cite{yarygina2018overcoming}:
\begin{itemize}
    \item \textbf{\texttt{polibio-gateway} (Puerto 8000):} Punto de entrada único que gestiona la autenticación perimetral, despacho estático de la interfaz SPA e inyección de encabezados de seguridad.
    \item \textbf{\texttt{polibio-doc-service} (Puerto 8002):} Motor documental oficial encargado de la formulación, correlativos y persistencia.
    \item \textbf{\texttt{polibio-crypto-service} (Puerto 8003):} Autoridad Certificadora Raíz (Root CA), motor PKI y ejecutor de cifrado de sobre.
    \item \textbf{\texttt{polibio-ledger-service} (Puerto 8004):} Motor de blockchain permisionado inmutable.
    \item \textbf{\texttt{polibio-ai-service} (Puerto 8005):} Orquestador de inferencia local acoplado a Ollama.
    \item \textbf{\texttt{polibio-telemetry-service} (Puerto 8006):} Colector de métricas en tiempo real integrado con Prometheus.
\end{itemize}

\subsection{Capa de Persistencia Híbrida y Row-Level Security (RLS)}
La persistencia combina un motor relacional PostgreSQL 16 con almacenamiento estático estructurado. Se diseñaron disparadores (\textit{triggers}) y políticas RLS donde cada consulta evalúa la variable de sesión:
\begin{equation}
\begin{split}
\text{RLS\_Filter}(row) \equiv (row.\text{level} \le \\
\text{current\_setting}('\text{app.clearance}'))
\end{split}
\end{equation}

\section{Seguridad Multinivel y Cifrado de Doble Sobre}
\subsection{Formalización Matemática del Modelo de Seguridad en Polibio}
Sea $\mathcal{U}$ el conjunto de funcionarios y $\mathcal{D}$ el conjunto de documentos oficiales. Se define la función de clasificación $\lambda: \mathcal{D} \to \{1, 2, 3, 4, 5\}$ y la función de habilitación $\mu: \mathcal{U} \to \{1, 2, 3, 4, 5\}$, donde:
\begin{align*}
1 &\equiv \text{Público}, & 2 &\equiv \text{Uso Interno}, \\
3 &\equiv \text{Reservado}, & 4 &\equiv \text{Confidencial}, \\
5 &\equiv \text{Secreto de Estado}.
\end{align*}

\subsection{Algoritmo de Cifrado y Ensamblaje del Sobre Criptográfico}
Para todo documento $d \in \mathcal{D}$ con $\lambda(d) \ge 3$, el microservicio criptográfico ejecuta el procedimiento formalizado en el Algoritmo 1.

\begin{lstlisting}[language=C, caption={Algoritmo de Sobre Criptográfico Híbrido (AES-256-GCM + RSA-OAEP)}]
// 1. Generacion de Clave de Cifrado de Contenido (CEK)
K_sym = CryptographicallySecureRandom(32); // 256 bits
IV    = CryptographicallySecureRandom(12); // 96 bits

// 2. Cifrado Autenticado AEAD
(C, Tag) = AES_256_GCM_Encrypt(Plaintext, K_sym, IV, AAD);

// 3. Envoltura Asimetrica para el Destinatario
WrappedKey = RSA_OAEP_Encrypt(K_sym, RecipientPublicKey);

// 4. Sanitizacion en Reposo
DatabaseBody = "[CONTENIDO CIFRADO - CLASIFICACION NIVEL " + Level + "]";
Envelope = { C, IV, Tag, WrappedKey, RecipientThumbprint };
\end{lstlisting}

\subsection{Descifrado al Vuelo en Memoria Volátil}
Cuando un usuario $u \in \mathcal{U}$ con $\mu(u) \ge \lambda(d)$ solicita el documento, el microservicio documental solicita la apertura del sobre a \texttt{polibio-crypto-service}. La clave de contenido $K_{sym}$ es recuperada en RAM, el texto $P$ es descifrado e inyectado directamente en el flujo de respuesta HTTP seguro (TLS), destruyendo las copias en memoria tras la serialización.

\section{Notarización en Blockchain y Auditoría Forense}
\subsection{Estructura del Bloque Criptográfico}
Cada bloque $B_k$ del libro mayor permisionado satisface la relación de encadenamiento:
\begin{equation}
H_k = \text{SHA-256}(k \parallel t_k \parallel H_{k-1} \parallel \text{MerkleRoot}(\mathcal{T}_k) \parallel nonce)
\end{equation}
donde $\mathcal{T}_k$ es el conjunto de transacciones validadas en el período temporal correspondiente.

\subsection{Trazabilidad Inmutable de Infracciones de Acceso (Forensic Ledger)}
A diferencia de los sistemas de auditoría en tablas SQL mutables, Polibio emite una transacción en la cadena ante cualquier intento fallido:
\begin{equation}
\begin{split}
\tau_{\text{violación}} = \langle \text{MLS\_ACCESS\_REJECTED}, d_{id}, \\
\mu(u), \lambda(d), u_{\text{thumbprint}}, t \rangle
\end{split}
\end{equation}
Esta transacción es incorporada al siguiente bloque, asegurando que ningún administrador o atacante pueda purgar la evidencia de una tentativa de espionaje o acceso ilícito.

\section{Inferencia Neuronal Soberana On-Premise}
\subsection{Aislamiento y Mitigación de Fuga de Telemetría}
El microservicio \texttt{polibio-ai-service} interactúa de forma exclusiva con una instancia local de Ollama que ejecuta modelos abiertos compactos (e.g., \textit{Qwen-2.5-0.5B}). El contenedor carece de interfaz de enrutamiento hacia Internet, asegurando que:
\begin{equation}
\mathcal{P}_{\text{exfiltración}} \equiv 0
\end{equation}

\subsection{Asistencia en Redacción Canónica y Formulación de Oficios}
El motor asiste a los redactores en la adecuación estilística y estructuración tripartita de oficios y memorandos, respetando las restricciones de correlativos predeterminadas por cada departamento institucional.

\section{Evaluación Experimental y Resultados}
\subsection{Entorno Experimental y Configuración de Pruebas}
Las evaluaciones se condujeron en un clúster contenerizado sobre Windows Subsystem for Linux (WSL2), con procesador x86-64 y memoria RAM de 16 GB. Se instrumentaron bancos de pruebas automatizados mediante los scripts de verificación \texttt{test\_confidentiality.js}, \texttt{test\_dept\_templates\_e2e.js} y \texttt{test\_image\_upload.js}.

\subsection{Latencia Criptográfica y Sobrecarga Computacional}
La Tabla I resume los tiempos medios observados para las operaciones criptográficas críticas:

\begin{table}[htbp]
\caption{Rendimiento de Operaciones Criptográficas en Polibio}
\begin{center}
\resizebox{\columnwidth}{!}{
\begin{tabular}{lcc}
\toprule
\textbf{Operación Criptográfica} & \textbf{Tiempo Medio (ms)} & \textbf{Desv. Est. (ms)} \\
\midrule
Generación Hash Canónico SHA-256 & 0.42 & 0.08 \\
Cifrado de Sobre AES-256-GCM      & 1.18 & 0.15 \\
Envoltura Asimétrica RSA-OAEP-2048 & 3.45 & 0.32 \\
Descifrado al Vuelo en Memoria   & 1.22 & 0.11 \\
Sellado y Minería de Bloque Ledger & 2.10 & 0.28 \\
Firma Digital PKI X.509          & 4.80 & 0.41 \\
\bottomrule
\end{tabular}
}
\end{center}
\end{table}

\subsection{Verificación de Cero Fuga y Resistencia Forense}
Las inspecciones directas a las imágenes de almacenamiento secundario y tablas de base de datos confirmaron que el 100\% de los documentos con clasificación $L \ge 3$ se encuentran totalmente ofuscados y protegidos con llaves AEAD. Asimismo, los intentos de violación de acceso generaron los registros criptográficos inalterables en el bloque correspondiente de la cadena.

\section{Conclusiones y Trabajo Futuro}
En este trabajo se ha presentado \textbf{Polibio}, una plataforma de microservicios soberana que solventa los vacíos de seguridad, soberanía y trazabilidad de los gestores documentales contemporáneos. Mediante la articulación armónica del modelo Bell-LaPadula, cifrado de sobre autenticado en reposo, una cadena de bloques permisionada con registro forense y un asistente de IA on-premise estricto, Polibio ofrece un entorno blindado para comunicaciones institucionales de alto valor estratégico.

El trabajo futuro contempla la incorporación de esquemas criptográficos poscuánticos (post-quantum cryptography, PQC) como ML-KEM y ML-DSA para la protección a largo plazo contra la amenaza cuántica, así como la federación de cadenas de bloques entre múltiples entes gubernamentales independientes.

\bibliographystyle{IEEEtran}
\bibliography{referencias}

\end{document}
